\chapter{Живые опыты: пульт, полёты и зуммер}\label{ch:pult}

До сих пор мы читали. В этой главе то, что мы \emph{делали} с живыми устройствами: управляли дроном через USB-пульт и заставили зуммер на модуле играть мелодию по Wi-Fi. Все факты проверены на оборудовании владельца (ноябрь 2026).

\section{USB-пульт Wing}

Пульт со стиками подключается к компьютеру USB-кабелем и определяется как адаптер CP210x (порт \texttt{COM3}, 115\,200 бод). LiteBee Go умеет управлять через него дроном Wing, наша программа тоже (команды \texttt{wing-status} и \texttt{wing-fly}). \pr{}

\begin{tabularx}{\textwidth}{@{}p{3.6cm} Y@{}}
\toprule
\hdr{Что} & \hdr{Как} \\ \midrule
Кадры & К пульту: \texttt{FD FD} длина данные. От пульта: \texttt{FC FC} длина данные. \\
Связь & Приветствие \texttt{FD FD 01 E9}, ответ \texttt{FC FC 02 E9 03}; затем примерно 20 раз в секунду пульт присылает пакет телеметрии (30 байт: положение, углы, батарея, режим). Нужен пульс: один байт \texttt{00} каждые 500 мс. \\
Режим пульта & Должен быть режим программирования (индикатор L4 горит зелёным, переключает кнопка K-0). В обычном режиме данные застывают. \\
Привязка дрона & Кнопка K-0 и правый стик в верхний правый угол (индикатор L3 загорается зелёным). \\
Сообщения & \texttt{vel\_control} 27, \texttt{pos\_control} 30, \texttt{set\_rotate} 18, зуммер 25, прямая передача 26. \\ \bottomrule
\end{tabularx}

\subsection*{Полёты}

Пять полётов Wing через пульт (дата и результаты в \texttt{docs/PROTOCOL\_SINGLE.md}). Первые три (октябрь) показали, что \textbf{команда сдвига \texttt{[30, 4, X, Y]} не работает}, а команда скорости \texttt{[27, 2, 0, Y, X, 0]} работает. В двух полётах 7 ноября дрон выполнил план без вмешательства: калибровка, взлёт, высота 80 см, шаг вперёд на 2 с, шаг вправо на 2 с, посадка. \pr{}

\begin{itemize}
  \item \textbf{Единицы скорости} --- см/с: при 15 см/с за 2 с дрон ушёл на 29--36 см. \pr{}
  \item \textbf{Выбег после «стоп»}: ещё 10--13 см. \pr{}
  \item Взлёт стандартной командой поднимает дрон почти на метр (90--98 см); команда высоты опускает до 80--86 см. \pr{}
  \item Положение в данных пульта живое: $x$ вперёд, $y$ вправо (сантиметры), точность около 1--3 см относительно заданного. \pr{}
\end{itemize}

\begin{opasno}[Полёты только осознанно]
Наша программа перед полётом проверяет связь, ровность дрона на земле, заряд батареи (не ниже 7,4 В) и отсутствие файла \texttt{LAND\_NOW}, а в полёте сажает дрон при наклоне больше 30\textdegree, высоте больше 130 см, потере данных на секунду, батарее ниже 6,6 В, любой ошибке или Ctrl+C. Полёт всегда запускается отдельным флагом \texttt{--go}, после того как владелец подтвердил, что площадка свободна, пропеллеры стоят, а пульт рядом.
\end{opasno}

Пульт не используется дронами формации (FM): они работают только по Wi-Fi с программой на компьютере. \vyv{}

\section{Зуммер на модуле}

К разъёму модуля с подписями \textbf{U0TX, U0RX, 5V, GND} подключён зуммер с четырьмя проводами. Это UART0 процессора модуля (\texttt{/dev/ttyS0}, 115\,200 бод, 8N1), свободный от связи с дроном. \pr{}

\subsection*{Протокол}

Зуммер понимает тот же вид кадров, что и модуль: \texttt{AD FA} команда данные сумма, сумма равна (команда плюс данные) по модулю 256. Нота играется командой \texttt{01}:

\begin{center}\ttfamily AD\ FA\ 01\ <нота>\ <длительность>\ <сумма>\end{center}

Нота взята из таблицы LiteBee Go: 0 это B0, до третьей октавы 25, ре 27, ми 29, фа 30, соль 32, ля 34, си 36, до четвёртой октавы 37. Длительность в единицах около 0,05 с (10 даёт примерно полсекунды). Пример: \texttt{AD FA 01 19 0A 24}. \pr{} (проверено на слух: гамма и «С днём рождения» звучат верно)

При включении модуля зуммер сам пискнул: штатная программа отправляет в этот порт несколько байт, и зуммер принял их за свою команду. \pr{}

\subsection*{Мелодия по Wi-Fi через дрон}

Формационный протокол звука не знает, поэтому путь идёт в обход:

\begin{figure}[tbp]
\centering
\resizebox{\textwidth}{!}{%
\begin{tikzpicture}[
  box/.style={draw=rule, fill=paper, rounded corners=3pt, minimum height=12mm, minimum width=28mm, align=center, font=\sffamily\scriptsize},
  hot/.style={box, draw=accent, fill=accentsoft},
  lnk/.style={-{Stealth[length=2.2mm]}, line width=0.8pt, color=accentdark},
  tag/.style={font=\sffamily\tiny, color=inkgray, align=center}
]
\node[box] (pc) {Компьютер\\\texttt{litebee buzz}};
\node[hot, right=17mm of pc] (rl) {Ретранслятор\\на дроне};
\node[box, right=17mm of rl] (md) {Модуль\\команда \texttt{FA}};
\node[box, right=17mm of md] (bz) {Зуммер};
\draw[lnk] (pc) -- node[above, tag]{UDP 5555\\Wi-Fi} (rl);
\draw[lnk] (rl) -- node[above, tag]{\texttt{ttyS0}\\дрона} (md);
\draw[lnk] (md) -- node[above, tag]{UART0\\модуля} (bz);
\end{tikzpicture}}
\caption{Путь ноты с компьютера до зуммера.}\label{fig:buzzpath}
\end{figure}

\begin{enumerate}
  \item \textbf{Кадр-реле.} Команда модуля \texttt{FA} пересылает байты на его UART0, то есть на зуммер. Кадр: \texttt{AD FA FA 08} и 8 байт (кадр ноты, добитый нулями) и сумма. Пример, который сыграл ноту: \texttt{AD FA FA 08 AD FA 01 19 0A 24 00 00 F1}. \pr{}
  \item \textbf{Ретранслятор на дроне.} На дроне нет ни Python, ни \texttt{nc}, ни компилятора, поэтому программу в 290 байт мы собрали вручную в машинном коде ARM (\texttt{tools/firmware/build\_udp\_tty\_relay.py}, проверено дизассемблером). Она слушает UDP-порт 5555 и пишет в \texttt{/dev/ttyS0} пакеты, начинающиеся с \texttt{AD FA}. \pr{}
  \item \textbf{Компьютер.} \texttt{python -m litebee buzz --host <IP дрона> --go "C3:10 E3:10 G3:20"}. Без \texttt{--go} печатается только план. Песня «С днём рождения» с компьютера по Wi-Fi прозвучала целиком. \pr{}
\end{enumerate}

\begin{vnimanie}[Ограничения]
\begin{itemize}
  \item Ретранслятор живёт \emph{в оперативной памяти дрона} и пока открыт сеанс ADB. После перезагрузки или отключения кабеля его нужно запускать заново. Автозапуск потребовал бы записи во флеш (\texttt{autorun.sh}); мы этого не делали. Переживает ли процесс отключение кабеля, не проверено. \nei{}
  \item Пока он работает, порт 5555 открыт для всей сети дрона (раздел <<Безопасность>> главы~\ref{ch:sys}).
  \item Громкость, паузы и остановку ноты мы не знаем; проверены ноты 25--44.
\end{itemize}
\end{vnimanie}

\section{Что из этого следует для шоу}

Один дрон формации может «петь», пока остальные летят: ноты уходят отдельными UDP-пакетами по расписанию, которое знает программа на компьютере. Для этого в приложение нужен шаг <<спеть>>, а на дроне работающий ретранслятор. Вживую вместе с полётом это мы не проверяли. \vyv{}
